\documentclass[11pt]{amsart}
\usepackage{calc,amssymb,amsmath,ulem,stmaryrd,fullpage}
\usepackage{enumerate}
\usepackage{alltt}
\RequirePackage[dvipsnames,usenames]{color}
\let\mffont=\sf
\normalem
%\input{kmacros3.sty}
\input{xy}
\xyoption{all}
\usepackage{hyperref}
\usepackage{graphicx}
\usepackage[all,cmtip]{xy}
\usepackage{verbatim}
\usepackage[left=0.95in,top=0.95in,right=0.95in,bottom=0.95in]{geometry}
\newcommand{\tauCohomology}{T}
\newcommand{\FCohomology}{S}


%\theoremstyle{theorem}
%\newtheorem*{mainthm*}{Main Theorem}

\newcommand{\note}[1]{\marginpar{\sffamily\tiny #1}}

\def\todo#1{\textcolor{Mahogany}%
{\sffamily\footnotesize\newline{\color{Mahogany}\fbox{\parbox{\textwidth-15pt}{#1}}}\newline}}

\renewcommand{\O}{\mathcal O}
\newcommand{\ie}{{\itshape i.e.} }
\newcommand{\roundup}[1]{\lceil #1 \rceil}
\newcommand{\rounddown}[1]{\lfloor #1 \rfloor}
\renewcommand{\to}[1][]{\xrightarrow{\ #1\ }}
\newcommand{\onto}[1][]{\protect{\xrightarrow{\ #1\ }\hspace{-0.8em}\rightarrow}}
\newcommand{\into}[1][]{\lhook \joinrel \xrightarrow{\ #1\ }}
%\newcommand{DBDual}[1]{{\underline{\omega}_{#1}}}
\newcommand{\DBDual}[1]{{{\uuline \omega}{}^{\mydot}_{#1}}}
\newcommand{\Tt}{{\mathfrak{T}}}
%\newcommand{\bn}{{\mathfrak{n}} }
\newcommand{\sol}[1]{ \vskip 6pt{\color{blue} {\bf Solution:} #1} \vskip 6pt}

\newcommand{\bR}{{\mathbb{R}}}
\newcommand{\va}{{\vec{a}}}
\newcommand{\vb}{{\vec{b}}}
\newcommand{\vzero}{{\vec{0}}}
\newcommand{\vi}{{\vec{i}}}
\newcommand{\vj}{{\vec{j}}}
\newcommand{\vk}{{\vec{k}}}
\newcommand{\vh}{{\vec{h}}}
\newcommand{\vr}{{\vec{r}}}
\newcommand{\vu}{{\vec{u}}}
\newcommand{\vv}{{\vec{v}}}
\newcommand{\vw}{{\vec{w}}}
\newcommand{\vx}{{\vec{x}}}
\newcommand{\vy}{{\vec{y}}}
\newcommand{\vz}{{\vec{z}}}
\newcommand{\vT}{{\vec{T}}}
\newcommand{\vN}{{\vec{N}}}
\newcommand{\vF}{{\vec{F}}}
\newcommand{\vG}{{\vec{G}}}
\newcommand{\bF}{\mathbb{F}}
\newcommand{\bQ}{\mathbb{Q}}
\newcommand{\bZ}{\mathbb{Z}}
\newcommand{\Inn}{\mathrm{Inn}}
\newcommand{\Aut}{\mathrm{Aut}}
\newcommand{\Hom}{\mathrm{Hom}}
\newcommand{\End}{\mathrm{End}}

\begin{document}
\title{HW \#4 -- Math 6310\\Fall 2019}
\author{Due: Friday, October 25th}
\maketitle
\noindent
{\bf 1.}  Show that the set $\bZ[\sqrt{2}] \subseteq \bR$ of real numbers of the form $m + n \sqrt{2}$ with $m, n \in \bZ$, is a Euclidean domain with respect to the function $\delta(m + n\sqrt{2}) = |m^2 - 2n^2|$.  
\newpage
\noindent
{\bf 2.}  Show that if $D$ is a PID and $0 \neq a \in D$, then $D/(a)$ is a field if $a$ is irreducible and $D/(a)$ is not an integral domain if $a$ is not irreducible.
\newpage
\noindent
{\bf 3.}  Prove the following irreducibility criterion of Eisenstein.  If $f(x) = a_0 + a_1 x + \dots + a_n x^n \in \bZ[x]$ and there exists a prime $p \in \bZ$ such that $p | a_i$ for all $i < n$ and $p^2 \not\,| a_0$, then $f(x)$ is irreducible in $\bQ[x]$.  
\vskip 3pt
\emph{Hint:}  Factor $f = g \cdot h$.  Consider what the leading terms and constant terms of $g$ and $h$.  Reduce them modulo $p$ as we did in the proof of Gauss' Lemma.
\newpage
\noindent
{\bf 4.}  Suppose $D$ is a commutative integral domain which is not a field.  Prove that $D[x]$ is not a PID.
\newpage
\noindent
{\bf 5.}  Suppose that $R$ is a commutative ring and $M$ is an $R$-module.  Suppose that $\phi : M \to R^{\oplus n}$ is surjective.  
Show that there exists a map $\psi : R^{\oplus n} \to M$ such that $\phi(\psi(x)) = x$.  
\vskip 3pt
Show further that $M$ is isomorphic to $R^{\oplus n} \oplus N$ for some other $R$-module $N$.
\vskip 3pt
\emph{Hint: }  Let $N$ be the kernel of $\phi$.  Write down a map from $N \oplus R^{\oplus n}$ to $M$ and show it is an isomorphism.
\newpage
\noindent
{\bf 6.}  Suppose that $p$ is prime.  Show that $\bZ/\langle p^e \rangle$ is not a direct sum of any two nonzero submodules.  Is the same true for $\bZ$?
\newpage
\noindent
{\bf 7.}  Determine $\Hom_{\bZ}(\bZ/\langle m \rangle, \bZ/\langle n \rangle)$ for all $m, n > 0$.
\newpage
\noindent
{\bf 8.}  A left $R$-module $M$ is called \emph{irreducible} if $M \neq 0$ and $0$ and $M$ are the only submodules of $M$.  Show that $0 \neq 0M$ is irreducible if $M$ is cyclic with ever non-zero element a generator.
\newpage
\noindent
{\bf 9.}  A left ideal $I \subseteq R$ is called \emph{maximal} if $R \neq I$ and there is no left ideal $I'$ with $I \subsetneq I' \subsetneq R$.  Show that $M$ is irreducible if and only if $M \cong R/I$ where $I$ is a maximal left ideal of $R$.
\newpage
\noindent
{\bf 10.}  Show that if $M_1$ and $M_2$ are irreducible left $R$-modules, then any nonzero homomorphism $M_1 \to M_2$ is an isomorphism.  Conclude that if $M$ is irreducible, then $\End_R(M) = \Hom_R(M, M)$ is a division ring.
\end{document}

